Cross-model correlations between alignment benchmarks are routinely interpreted as evidence of alignment coherence, but this interpretation conflates capability-driven co-movement with genuine alignment-specific structure. We show that controlling for MMLU --- a proxy for general language model capability --- reduces the Spearman correlation between TruthfulQA MC2 (factuality) and a bias proxy across $N{=}296$ open-weight LLMs from 0.732 to 0.343, a 53\% reduction that is statistically decisive (Fisher $z = 6.97$, $p = 3.22 \times 10^{-12}$). MMLU alone explains 49\% of TruthfulQA variance and 76\% of the bias-proxy variance, confirming scale confounding of substantial magnitude. A significant residual partial correlation (0.343, $p = 1.40 \times 10^{-9}$) indicates genuine alignment-specific co-movement beyond scale, though its structural interpretation remains ambiguous (BCa CI [0.180, 0.492] spans a pre-specified scenario threshold at $N{=}296$). We propose the Fisher $z$ raw-vs-partial difference test as a practical diagnostic for scale confounding in alignment benchmark correlation analysis, and demonstrate it via a four-step pre-registered verification pipeline with 57/57 test cases passing. Our findings suggest that positive cross-model alignment benchmark correlations substantially overestimate alignment-specific coherence, and that capability control should be standard practice before interpreting benchmark co-movement as evidence of aligned training objectives.
